\section{Discussion and limitations}
\label{sec:discussion}
\paragraph{Lifecycle coverage and delivery reliability.}
DRAFT-Bench provides evidence across requirement discovery, creative exploration, design and revision, and production handoff. Its Requirement-to-Artifact framework makes the linkage between evolving obligations and delivered versions explicit. The three findings separate stage-specific capability, requirement compliance, and the interpretation of missing or ambiguous evidence. Astra completes more of the observed execution workload, while Claude has a higher observed weighted contribution on the selected shared Brief Bank core. Missing outcomes, unknown rubric contributions, and different tools prevent a total quality ranking. Neither routing and planning scores nor completed execution establish independent professional acceptance.

\paragraph{Intake coverage is a separate capability.}
RFQ's near-complete average field coverage does not establish that every critical output requirement was elicited, that the final brief is usable, or that a downstream design will be accepted. The success ordering changes between the original and additional cohorts and narrows under the exclusion sensitivity check. Auxiliary models, continuation history, and scripted disclosure also change, so this is descriptive heterogeneity rather than an isolated domain effect. A claim that either system is generally better at elicitation exceeds these data.

\paragraph{The evaluator is part of the measurement.}
The unchanged-artifact R1 pilot demonstrates sensitivity to acceptance rules and evidence presentation: a text-ranking task need not require an image, and a missing edit may be an unfulfilled obligation, justified clarification, or an invalid feedback prerequisite. These states require different labels. R1's remaining variation and lack of expert adjudication prevent calling the new scores corrected ground truth. Its seven selected pairs cannot estimate a benchmark-wide grading error rate. The main scorer's image sampling and incomplete dialogue view add limits: evidence available in the archive is not necessarily evidence delivered to the judge.

\paragraph{Implications for agent and evaluator development.}
Track active obligations and approved properties across versions; make the current checkpoint and the eventual deliverable set explicit. Reconcile brief-level scope with turn-level delta instructions before judging an omission. Provide artifact facts independently of candidate claims, and retain the candidate's outward explanation as evidence of communication, not as proof that a visual edit occurred. Separately label delivery completeness, the reason for blocking, and whether escalation was justified. These are development targets motivated by the audit, not interventions validated by this study. They align with professional-workflow evaluation~\citep{openai2026ironclad} and guidance on outcomes and grader calibration~\citep{anthropic2026evals}.

\paragraph{The next discriminating experiments.}
First use disputed cases to specify acceptance rules, then freeze those rules and evaluate a separate, stratified sample with experts blinded to system identity, including flagged and unflagged outputs and blocker decisions. Second compare judges on the same frozen evidence and report agreement and errors against expert labels; keep the R1 development cases separate from that test. Third compare agents with matched tools and repeated runs, and connect vague, elicited, and complete briefs to one downstream executor. This would test RFQ's delivery value and enable validation of an integrated RFQ-to-handoff chain. Production experts should separately assess plate placement, preflight, and factory acceptance before making design-to-manufacturing claims. None of these experiments is completed here.

\paragraph{Scope and causal limits.}
The 272 lineages and 301 protocol units span curated, synthetic, and historical tasks; the 333 records include duplicate associations and are not independent observations. Design executions are single included trajectories. RFQ variants share source families, its original amendment is post-hoc, and the cohorts use different auxiliary models and attempt histories. Timeouts, infrastructure exclusions, and missing quality dimensions remain distinct. Tools, runtimes, platform instructions, isolation, and dates differ across systems, preventing model-weight attribution or a causal quality--reliability trade-off. Human acceptance calibration remains absent. Production handoff is evaluated through artifacts and file evidence; hashes, pixel differences, and a White separation resource cannot certify professional quality or factory acceptance. Some rejected historical judge outputs were discarded.

\section{Conclusion}
DRAFT-Bench evaluates frontier agents across the professional design lifecycle through separate task protocols. Its central contribution is Requirement-to-Artifact evaluation: linking evolving requirements, protected properties, artifact versions, and delivery obligations to inspectable evidence. Across 35 design briefs and 43 conditions, Astra completes 43/43 and Claude 35/43, while full structural credit coexists with semantic flags and unresolved deliveries. The separate 40-scenario RFQ track shows high coverage and cohort-sensitive success ordering; the shared-case totals are 36/39 for Astra and 34/39 for Claude. SKU expansion, version dependencies, and White ink files show why human decisions and production constraints must be part of evaluating delivery. These results support a stage-specific capability analysis, with distinct measures of execution, requirement compliance, and justified escalation. They do not establish a pure model ranking, expert-accepted delivery, or factory acceptance. A paired RFQ-to-handoff experiment and independent expert calibration are the next steps beyond the current lifecycle-wide benchmark.

\paragraph{Data and code availability.}
Public materials are at \url{https://github.com/super-inbox/DRAFT-Bench}. Data release \texttt{draftbench-v0.1.0} (CC BY 4.0) contains the briefs, 24 input fixtures, RFQ scenarios and 100 run records with offline score replay, numeric records for all 301 protocol units, and the R1 scores. Toolkit release \texttt{draftbench-v0.2.0} (new code under MIT) runs these tasks under a new protocol; it does not reproduce the historical environments. Task inputs for the auxiliary protocols, candidate outputs beyond our figures, complete traces, judge payloads, the semantic-check audit, case evidence, and this paper's analysis scripts are not public, so the design judgments cannot be re-created from these materials alone.
